\chapter{The Native Apparatus Presents The Residue}

\section{Forcing the residue to show}

The corridor has been turned and its faces mapped; the price of the turn has been named and its magnitude withheld. This
part brings the residue to a bench and forces it to show. Until now the residue has been carried, read, bracketed --- but
always as a quantity the construction tracked internally. Here it is put under an \emph{apparatus}, a physical arrangement
built to make the residue present itself as a definite reading, and the apparatus is the physical gauge's own: force, orbit,
field. This section says what it means to force the residue to show, and why the showing is the same invariant every earlier
chapter carried.

To force the residue to show is to arrange a loop whose failure to close the apparatus can read directly. Every bench in the
next section is such an arrangement: a torsion fiber whose twist reads a force, an orbit whose precession reads a binding, a
circuit whose holonomy reads a phase. The apparatus does not create the residue; it exposes it, by building a loop tight
enough that the residue's alteration lands on a dial. This is the experimental content of everything the book has said
abstractly: a measurement is a loop's alteration, and an apparatus is a loop engineered so the alteration is legible. The
residue was always there; the bench is what makes it show.

The important discipline of the part is fidelity to the benches that actually exist. The instrument carries a definite,
finite roster of apparatus, and the physical gauge reads exactly those and no others. It would be easy, and dishonest, to
write a bench for every experiment physics has ever run --- to add a spin-splitting bench here, a scattering bench there,
because the physics permits them. The construction does not carry them, and the book will not claim them. What it carries is a
short list --- a torsion balance, a bound orbit, a force law with its bracket structure, an electron equation, a variational
projection, a field-expulsion effect, a threshold effect --- and the gauge reads the residue on those, marking each reading
as a reading. An apparatus the construction does not hold is an apparatus the book does not describe. The roster is finite,
and fidelity to its finiteness is part of the honesty.

There is a reason the same residue can show on such different benches, and it is the unity Part III established. A torsion
balance and an electron equation look nothing alike, but each reads a contraction or a loop of the one curvature: the balance
reads its trace as a force, the equation reads its winding as a spinor. Because the faces are handlings of a single
invariant, a bench built to read any one handling is reading the same residue as a bench built to read another. This is why
the several benches of the next section do not measure several different things. They measure one residue, each forcing a
different face of it to show, and their agreement --- when it comes --- is the agreement of one invariant read many ways.

In the two verbs, forcing the residue to show is building a tange the apparatus can read off a funge the world holds. The
residue is the funge, the covariant invariant carried all along; the apparatus is a tange, a contravariant arrangement built
to select one characteristic and land it on a dial. To force the residue to show is to engineer the selection so the world's
bag deposits its reading where a bench can see it. The next section walks the roster of such tanges --- the physical benches
the construction actually carries --- and reads the residue on each.

\section{The physical bench}

This is the physical gauge's own section: the roster of benches the construction carries, each forcing the one residue to
show as a definite physical reading. The list is finite and fixed, and the section walks it in order --- a torsion balance,
a bound orbit, a force law with its bracket, an electron equation, a variational projection, a field-expulsion effect, a
threshold effect. Each bench reads a face of the residue; each reading is marked as a reading; and none is inflated past what
its apparatus delivers.

Begin with the torsion balance --- Cavendish's arrangement, a mass hung on a fiber whose twist reads the faintest
force~\cite{cavendish1798}. Here the residue shows as \emph{force}: the trace of the curvature, the Ricci part tied to the
local source, read off as the deflection of the fiber. The balance is a loop --- the fiber twists and the residue is the
angle it fails to return to zero --- and the reading is a torque, bracketed at the fiber's resolution. What the bench
delivers is the residue's trace face, the near field a local mass carries with it, made legible as a twist. The construction
supplies the trace; the identification of it with Newton's gravitational force is the physical reading, and Cavendish's fiber
is where it is cashed.

Next the bound orbit --- a body circulating in a central field, the residue shown as the \emph{orbit} and its
precession~\cite{newton1687,goldstein2002}. A closed orbit is a loop by construction, and the residue is exactly its failure
to close: the perihelion precession, the angle by which the orbit fails to return on itself after one circuit. This is
Chapter 5's slip made into an astronomical reading --- the loop that does not quite close, its anholonomy read as a
precession. The bench delivers the residue as the geometry of a binding, and the reading is an angle per orbit, bracketed at
the resolution of the tracking. The orbit does not measure a new residue; it reads the same curvature as the balance, through
a loop drawn in the sky instead of on a fiber.

The force law itself carries a bracket structure, and the third bench reads it. The Lorentz force $\frac{\mathrm{d}p_\mu}
{\mathrm{d}\tau} = e\,F_{\mu\nu}u^\nu$ writes the residue $F_{\mu\nu}$ as the operator that turns a state's advance into a
force~\cite{lorentz1904}, and the antisymmetric bracket structure of the dynamics --- the symplectic form that pairs a
coordinate with its conjugate --- is the arena the residue acts in. This is the residue shown as the \emph{coupling of a
charge to a field}: the field strength contracted against the state's four-velocity, delivering the force that bends the
worldline. The bench reads $F_{\mu\nu}$ directly, as the thing that couples charge to advance, and the bracket that pairs
position with momentum is the frame the reading is taken in. The residue here is the field strength doing its native job.

The fourth bench is the electron equation, Dirac's, where the residue shows as a \emph{spinor}~\cite{dirac1928}. The
apparatus is the demand for a first-order relativistic equation, and its output is the four-component object whose winding is
the electron of Chapter 8. Here the residue presents its loop face --- the double-cover winding, the phase that takes two
turns to close --- as a definite equation with a definite spectrum. The bench delivers the residue as the named, stable
reading it settled into: the electron, with its integer charge and half-integer spin, read off the structure the algebra
forced. This is the loop face of the curvature, exhibited as the particle the counting named.

The fifth bench is the variational projection --- Galerkin's method, where the residue shows as, precisely, a
\emph{residual}~\cite{galerkin1915}. Pose the field's equation in weak form, project it onto a finite family of test
functions, and demand that the leftover --- the residual $r = L\phi - f$ --- be orthogonal to every test function, $\langle
r, v_k \rangle = 0$. What the projection cannot drive to zero, the part of the residual orthogonal to the finite family, is
the residue in the plainest possible sense: the leftover the finite apparatus cannot resolve away. This bench is the book's
own thesis exhibited as a method --- a finite projection, an honest leftover it cannot bag, a residual carried rather than
discarded. The Galerkin residual is the resolution floor made into a numerical fact: what the finite basis cannot reach.

The sixth bench is field expulsion --- the Meissner arrangement and the proximity effect, where the residue shows as the
\emph{magnetic side}~\cite{meissner1933}. A superconductor expels an applied field, $B = 0$ inside, screening it over a
finite penetration depth; the proximity effect leaks that screening a finite distance into an adjacent material. Here the
residue presents its flux face --- the magnetic side of the field --- as a definite, finite screening length, a ratio of
applied field to shielding response. The bench delivers the field's other half, the magnetic side that must be present before
the field can close, and it delivers it finite: a penetration depth, not a continuum. The magnetic residue is read as the
depth to which the field is expelled.

The seventh bench is the threshold effect --- the photoelectric arrangement, where the residue shows as a \emph{quantum of
the coupling}~\cite{planck1901,planck1900}. Light below a threshold frequency ejects no charge however intense; above it,
each quantum $h\nu$ ejects one, $E = h\nu - W$. The bench reads the residue as the discrete threshold --- the quantized step
that turns a continuous input into a counted output --- and the reading is a count of quanta, not a smear of intensity. This
is the residue shown as the graininess of the coupling, the fact that the field meets matter in counted units. The threshold
is the residue's discreteness made into a bench a lamp and a plate can run.

In the two verbs, the roster is seven tanges over one funge. The residue is the single covariant invariant; each bench is a
contravariant apparatus tanging one characteristic out of it --- the balance the trace, the orbit the precession, the force
law the coupling, the equation the winding, the projection the residual, the expulsion the flux, the threshold the quantum.
Seven selections, one bag. The benches do not disagree, because they read one funge; they differ only in which tange they
perform. So the physical gauge reads the residue seven honest ways, each marked, each finite, each floored --- and the next
section names the sibling bench this volume does not own, where the same residue is read without any lab hardware at all.

\section{The other bench}

The physical benches read the residue as force, orbit, field, spinor, residual, flux, and quantum --- everything an apparatus
of matter can make show. This closing section names, briefly and then hands off, the sibling bench the same residue answers
to without any matter at all: a bench of pure elaboration, where the residue is read off the cost of constructing the reading
rather than off a fiber or a plate. It is not this volume's to describe, and the section's honesty is in naming it as another
gauge's and stopping there.

The claim that founds the other bench is that the residue shows not only when matter meets it but when the construction
merely \emph{makes} the reading. There is a cost to elaborating any reading --- an effort the instrument expends to carry the
residue to where it can be seen --- and that effort is itself a quantity, readable, bracketed, floored, exactly as a force or
a flux is. The sibling gauge builds its apparatus out of that effort: it weighs the construction of the reading and finds the
residue in the weighing. Where the physical bench forces the residue to show by arranging matter into a loop, the sibling
bench forces it to show by arranging the construction into a loop and reading what the construction is charged. Same residue,
no hardware.

This volume does not read that bench, and the reason is the partition of Chapter 7. The sibling gauge performs a projection
this volume does not: it reads the effort-face of the residue, the channel orthogonal to the force-face and the field-face
this volume owns. A projector sees only its own face; the physical gauge is blind to the effort-face exactly as it is blind
to the other gauges' readings, not from incompleteness but from orthogonality. So the book names the other bench, marks that
the same residue is read there, and hands it to the volume whose projection it is. What matters here is only that the residue
answers to a bench of pure construction as surely as to a bench of matter --- because that is the fact the whole four-volume
instrument will lean on at the end, when the benches that share no hardware turn out to close on one bracket.

There is a reason to name the other bench at all rather than pass it over, and it is the same reason the partition was a
strength. The physical bench and the elaboration bench share no apparatus: one is fibers and orbits and plates, the other is
pure construction with no matter in it. Two benches that share no hardware cannot be tuned to agree by any common adjustment,
because there is no common adjustment to make. So their eventual agreement on the coupling's bracket will be meaningful in a
way no single bench could earn. Naming the other bench now is planting the decorrelation the ending will harvest: the residue
shows on a bench of matter and on a bench of construction, and the two cannot see each other.

In the two verbs, the other bench is a tange this volume cannot perform on a funge it shares. The residue is the common
funge, read by every gauge; the effort-face is a tange the sibling gauge selects and this one cannot, because its projector
points the other way. So the book funges the residue whole, tanges off the faces it owns, and names the face it does not ---
handing the un-owned tange to the gauge whose projection reaches it. Part IV's first chapter thus closes with the residue
forced to show on the physical bench, seven ways, and the sibling bench named and handed off. The next chapter takes the
physical apparatus and turns it on its own act --- lets the bench weigh the cost of its own measuring --- and reads that
self-weight as a self-energy.
